% !TEX program = pdflatex
% Propuesta de rediseño UX para los módulos educativos de EduFEM
% Tesis - Software educativo de Elementos Finitos 2D (TP / DP)

\documentclass[11pt,a4paper]{article}

\usepackage[utf8]{inputenc}
\usepackage[T1]{fontenc}
\usepackage{lmodern}
\usepackage[spanish,es-noshorthands]{babel}
\usepackage[margin=2.2cm,headheight=14pt]{geometry}
\usepackage[protrusion=true,expansion=false]{microtype}
\usepackage{xcolor}
\usepackage{ifthen}
\usepackage{tikz}
\usepackage{tcolorbox}
\usepackage{amsmath,amssymb}
\usepackage{enumitem}
\usepackage{titlesec}
\usepackage{fancyhdr}
\usepackage{hyperref}
\usepackage{graphicx}
\usepackage{array}
\usepackage{booktabs}
\usepackage{caption}

\usetikzlibrary{shapes.geometric, shapes.misc, arrows.meta, positioning, fit,
                backgrounds, calc, decorations.pathreplacing, patterns,
                shadows.blur, shadings, 3d}
\tcbuselibrary{skins, breakable}

% --------- Paleta alineada con el tema darkly de la app ---------
\definecolor{bgDark}{HTML}{212529}
\definecolor{panelBg}{HTML}{2B2D33}
\definecolor{accent}{HTML}{4FA3FF}
\definecolor{phasePre}{HTML}{0D6EFD}
\definecolor{phaseProc}{HTML}{FD7E14}
\definecolor{phasePost}{HTML}{198754}
\definecolor{warn}{HTML}{FFC107}
\definecolor{danger}{HTML}{E74C3C}
\definecolor{good}{HTML}{43A047}
\definecolor{nodeOrange}{HTML}{FF9F43}
\definecolor{nodeBlue}{HTML}{4FC3F7}
\definecolor{midBlue}{HTML}{6FB8FF}
\definecolor{centerViolet}{HTML}{B86FFF}
\definecolor{selYellow}{HTML}{FFD54F}
\definecolor{softGray}{HTML}{45475A}
\definecolor{textMuted}{HTML}{8A8E96}

% --------- Estilos generales ---------
\hypersetup{
    colorlinks=true,
    linkcolor=accent,
    urlcolor=accent,
    citecolor=accent,
    pdftitle={Propuesta UX - Modulos Educativos EduFEM},
    pdfauthor={Tesis EduFEM}
}

\titleformat{\section}
  {\Large\bfseries\color{accent}}
  {\thesection}{0.7em}{}[{\color{softGray}\titlerule[0.6pt]}]

\titleformat{\subsection}
  {\large\bfseries\color{phaseProc}}
  {\thesubsection}{0.6em}{}

\titleformat{\subsubsection}
  {\normalsize\bfseries\color{phasePre}}
  {\thesubsubsection}{0.5em}{}

\pagestyle{fancy}
\fancyhf{}
\fancyhead[L]{\small\color{textMuted}EduFEM \textbar{} Propuesta UX}
\fancyhead[R]{\small\color{textMuted}\thepage}
\renewcommand{\headrulewidth}{0.4pt}
\renewcommand{\headrule}{\color{softGray}\hrule height \headrulewidth}

% --------- Macros para cajas y mockups ---------
\newtcolorbox{ideaBox}[1][]{
  enhanced,
  colback=panelBg!8,
  colframe=accent,
  boxrule=0.6pt,
  arc=2mm,
  left=8pt,right=8pt,top=6pt,bottom=6pt,
  fonttitle=\bfseries\color{white},
  coltitle=white,
  colbacktitle=accent,
  attach boxed title to top left={xshift=6mm,yshift=-3mm},
  boxed title style={arc=1.5mm,boxrule=0pt},
  title={#1}
}

\newtcolorbox{warnBox}[1][]{
  enhanced,
  colback=warn!8,
  colframe=warn!80!black,
  boxrule=0.6pt,
  arc=2mm,
  left=8pt,right=8pt,top=6pt,bottom=6pt,
  title={#1}
}

\newtcolorbox{moduleBox}[2][]{
  enhanced,
  colback=panelBg!5,
  colframe=#2,
  boxrule=1pt,
  arc=3mm,
  left=10pt,right=10pt,top=8pt,bottom=8pt,
  fonttitle=\large\bfseries\color{white},
  coltitle=white,
  colbacktitle=#2,
  attach boxed title to top left={xshift=10mm,yshift=-3.5mm},
  boxed title style={arc=2mm,boxrule=0pt},
  title={#1},
  breakable
}

% --------- Helpers TikZ para mockups de ventana ---------
\tikzset{
    win/.style={
        draw=softGray, fill=bgDark, rounded corners=2pt,
        line width=0.6pt
    },
    titlebar/.style={
        draw=softGray, fill=panelBg, rounded corners=2pt,
        line width=0.4pt
    },
    panel/.style={
        draw=softGray!60, fill=panelBg, rounded corners=1.2pt,
        line width=0.3pt
    },
    btn/.style={
        draw=accent, fill=accent!18, rounded corners=1pt,
        line width=0.3pt, text=white, font=\scriptsize
    },
    btnPhase/.style={
        draw=phaseProc, fill=phaseProc!22, rounded corners=1pt,
        line width=0.3pt, text=white, font=\scriptsize
    },
    chip/.style={
        draw=accent, fill=accent!30, rounded corners=2pt,
        line width=0.3pt, text=white, font=\tiny,
        inner sep=2pt
    },
    nodeDot/.style={circle, fill=nodeOrange, draw=white, line width=0.3pt,
                    minimum size=4pt, inner sep=0pt},
    midDot/.style={circle, fill=midBlue, draw=white, line width=0.3pt,
                    minimum size=2.8pt, inner sep=0pt},
    cnDot/.style={circle, fill=centerViolet, draw=white, line width=0.3pt,
                    minimum size=2.8pt, inner sep=0pt},
    selRing/.style={circle, draw=selYellow, line width=1.0pt,
                    minimum size=8pt, inner sep=0pt},
    label/.style={font=\tiny\color{white}, inner sep=1pt}
}

% ============================================================
\begin{document}

% --------- Portada ---------
\begin{titlepage}
\begin{tikzpicture}[remember picture, overlay]
  \fill[bgDark] (current page.south west) rectangle (current page.north east);

  % barra accent superior
  \fill[accent] ([yshift=-2.3cm]current page.north west) rectangle
    ([yshift=-2.5cm]current page.north east);

  % barras de fase
  \fill[phasePre]  ([yshift=-2.5cm]current page.north west) rectangle
    ([xshift=7cm,yshift=-2.65cm]current page.north west);
  \fill[phaseProc] ([xshift=7cm,yshift=-2.5cm]current page.north west) rectangle
    ([xshift=14cm,yshift=-2.65cm]current page.north west);
  \fill[phasePost] ([xshift=14cm,yshift=-2.5cm]current page.north west) rectangle
    ([yshift=-2.65cm]current page.north east);

  \node[anchor=north west, text=white, font=\fontsize{30}{34}\bfseries]
    at ([xshift=2.2cm,yshift=-3.8cm]current page.north west)
    {Rediseño UX};
  \node[anchor=north west, text=accent, font=\fontsize{22}{26}\bfseries]
    at ([xshift=2.2cm,yshift=-5.2cm]current page.north west)
    {Módulos Educativos};
  \node[anchor=north west, text=textMuted, font=\Large]
    at ([xshift=2.2cm,yshift=-6.2cm]current page.north west)
    {EduFEM \textbar{} Elementos Finitos 2D \textbar{} TP / DP};

  % Tarjeta de mockup esquemático
  \begin{scope}[shift={([xshift=2.2cm,yshift=-9cm]current page.north west)}]
    \draw[win] (0,0) rectangle (16,7);
    \fill[panelBg, rounded corners=2pt] (0.05, 6.4) rectangle (15.95, 6.95);
    \node[anchor=west, text=white, font=\small\bfseries]
      at (0.3, 6.67) {EduFEM};
    \fill[danger,  rounded corners=1pt] (15.0, 6.55) rectangle (15.4, 6.85);
    \fill[warn,    rounded corners=1pt] (14.5, 6.55) rectangle (14.9, 6.85);
    \fill[good,    rounded corners=1pt] (14.0, 6.55) rectangle (14.4, 6.85);

    % left: 3 phase tabs
    \draw[panel] (0.3, 0.5) rectangle (5.5, 6.3);
    \fill[phasePre, rounded corners=1pt] (0.5, 5.8) rectangle (1.9, 6.15);
    \fill[phaseProc, rounded corners=1pt] (2.0, 5.8) rectangle (3.4, 6.15);
    \fill[phasePost, rounded corners=1pt] (3.5, 5.8) rectangle (4.9, 6.15);
    \node[text=white, font=\tiny] at (1.2, 5.97) {Pre};
    \node[text=white, font=\tiny] at (2.7, 5.97) {Proc};
    \node[text=white, font=\tiny] at (4.2, 5.97) {Post};

    \node[text=textMuted, font=\tiny, anchor=west] at (0.5, 5.4) {Tablas};
    \node[text=textMuted, font=\tiny, anchor=west] at (0.5, 5.0) {Educación};
    \node[text=textMuted, font=\tiny, anchor=west] at (0.5, 4.6) {Selección};

    % right: shared mesh canvas
    \draw[panel] (5.7, 0.5) rectangle (15.7, 6.3);
    \node[text=textMuted, font=\tiny, anchor=south west]
      at (5.85, 6.05) {MeshCanvas compartido};

    % sample mesh inside canvas
    \begin{scope}[shift={(7.0, 1.5)}, scale=0.55]
      \foreach \x in {0,1,2,3} {
        \foreach \y in {0,1,2} {
          \draw[softGray, line width=0.4pt]
            (\x,\y) rectangle ({\x+1},{\y+1});
          \node[nodeDot] at (\x,\y) {};
        }
      }
      \foreach \x in {0,1,2,3,4} {
        \node[nodeDot] at (\x,3) {};
      }
      \foreach \y in {0,1,2,3} {
        \node[nodeDot] at (4,\y) {};
      }
      % highlighted element
      \fill[selYellow, opacity=0.18] (1,1) rectangle (2,2);
      \draw[selYellow, line width=1pt] (1,1) rectangle (2,2);
    \end{scope}

    % overlay pill (educational module)
    \draw[panel, fill=bgDark, line width=0.5pt, rounded corners=2pt]
      (10.5, 4.1) rectangle (15.4, 5.9);
    \node[text=accent, font=\tiny\bfseries, anchor=north west]
      at (10.65, 5.85) {M2 - Matriz B};
    \node[text=textMuted, font=\tiny, anchor=north west]
      at (10.65, 5.55) {evaluacion en pg2};
    \draw[selYellow, line width=0.6pt] (10.7, 4.4) rectangle (15.25, 5.3);
  \end{scope}

  \node[anchor=south west, text=textMuted, font=\small]
    at ([xshift=2.2cm,yshift=2cm]current page.south west)
    {Documento técnico de propuesta - Sin compromiso de implementación};
  \node[anchor=south west, text=accent, font=\small\bfseries]
    at ([xshift=2.2cm,yshift=1.5cm]current page.south west)
    {Tesis - Software Educativo FEM};
\end{tikzpicture}
\end{titlepage}

% --------- Resumen ejecutivo ---------
\section*{Resumen ejecutivo}
\addcontentsline{toc}{section}{Resumen ejecutivo}

Este documento propone un \textbf{rediseño UX diferenciado} de los nueve módulos
educativos de EduFEM (M0--M9). El diagnóstico parte de un patrón observado:
los nueve módulos heredan el mismo layout fijo (header + sidebar 280\,px +
visualización + footer), lo que genera una sensación de uniformidad que
\emph{contradice} la diversidad de los conceptos enseñados. Un módulo sobre
calidad geométrica de la malla y otro sobre el círculo de Mohr no deberían
verse iguales: cada concepto tiene una geometría natural de presentación.

La propuesta articula tres ejes complementarios:

\begin{enumerate}[leftmargin=1.2em, itemsep=2pt]
  \item \textbf{Tres modos de presentación} en lugar de uno solo:
        \emph{Toplevel autónomo}, \emph{Overlay sobre el canvas compartido} y
        \emph{Embed en el panel lateral de la fase}. Cada módulo elige el modo
        que su concepto pide.
  \item \textbf{Operación desde el canvas principal}: el elemento bajo análisis,
        el nodo seleccionado y la iluminación se manipulan desde el
        \texttt{MeshCanvas} compartido y la barra lateral de la fase activa.
        Los módulos consumen esa selección, no la duplican.
  \item \textbf{Layouts diferenciados}: cada módulo recibe una propuesta visual
        única, justificada pedagógicamente y mockup-eada con TikZ. Diez
        layouts, no uno repetido diez veces.
\end{enumerate}

\begin{ideaBox}[Premisas innegociables]
\textbf{Cero ruido visual estructural.} Todo widget debe justificarse contra
un objetivo pedagógico concreto. \textbf{Cero duplicación con el menú principal:}
si el dato (Q4/Q9, TP/DP, material, unidades) ya vive en \texttt{Modelo\,$\triangleright$\,...}
no se repite dentro del módulo. \textbf{Self-explaining.} El módulo debe
contar su historia con una animación al abrir y micro-tooltips contextuales,
no con bloques de texto en sidebar.
\end{ideaBox}

\tableofcontents
\newpage

% ============================================================
\section{Diagnóstico: por qué los módulos se sienten clonados}

\subsection{El layout único como cuello de botella}

\texttt{BaseEducationalModule} impone una caja con cuatro zonas fijas:

\begin{center}
\begin{tikzpicture}[scale=0.9]
  \draw[win] (0,0) rectangle (12, 7);
  % header
  \draw[panel] (0.1, 6.1) rectangle (11.9, 6.9);
  \node[text=white, font=\small, anchor=west] at (0.3, 6.5) {Header (titulo, badge, botones)};
  % sidebar
  \draw[panel] (0.1, 0.6) rectangle (3.0, 6.0);
  \node[text=white, font=\small, anchor=north, align=center, text width=2.4cm]
    at (1.55, 5.8) {Sidebar 280\,px\\(controles)};
  % viz
  \draw[panel] (3.1, 0.6) rectangle (11.9, 6.0);
  \node[text=white, font=\small] at (7.5, 3.3) {Visualizacion};
  % footer
  \draw[panel] (0.1, 0.1) rectangle (11.9, 0.5);
  \node[text=white, font=\small, anchor=west] at (0.3, 0.3) {Footer (animacion)};
  \node[text=textMuted, font=\tiny, anchor=south] at (6, -0.5)
    {Fig. 1 - Layout actual de \texttt{BaseEducationalModule}};
\end{tikzpicture}
\end{center}
\vspace{0.6cm}

El sidebar de 280\,px se convierte en un \emph{vertedero} de elementos:
combos de selección, sliders, párrafos explicativos, tarjetas de métricas
duplicadas, recordatorios teóricos. Esto genera tres patologías:

\begin{itemize}[leftmargin=1.2em, itemsep=2pt]
  \item \textbf{Sobre-carga cognitiva.} El alumno debe saltar entre
        ``leer texto del sidebar'' y ``mirar el gráfico'' para entender qué hace
        cada control. Los párrafos didácticos deberían ser \emph{ambientales}
        (overlay, tooltip al hover), no estructurales.
  \item \textbf{Duplicación con el modelo.} Selectores de elemento, valores
        de E, $\nu$, $t$, tipo de análisis aparecen en módulos cuando ya
        están definidos en \texttt{Modelo}. Cualquier desincronización
        confunde al alumno.
  \item \textbf{Layout incongruente con el concepto.} El círculo de Mohr
        (M8) merece pantalla con dos paneles vinculados; la matriz B (M2)
        merece la matriz LaTeX como protagonista. El sidebar fijo le quita
        protagonismo a ambas.
\end{itemize}

\subsection{Inventario de redundancias detectadas}

\begin{center}
\renewcommand{\arraystretch}{1.2}
\begin{tabular}{@{}l l l@{}}
\toprule
\textbf{Control} & \textbf{Módulos donde aparece} & \textbf{Fuente real} \\
\midrule
Selector Q4 / Q9 & M1 & \texttt{Modelo\,$\triangleright$\,Tipo de Elemento} \\
Selector TP / DP & (eliminado, ya OK) & \texttt{Modelo\,$\triangleright$\,Tipo de Análisis} \\
$E$, $\nu$ editables & M4 & Material del elemento \\
Espesor $t$ & M4 & Atributo del elemento \\
Combo de elemento & M0, M3, M4, M5 & Click en MeshCanvas \\
Tarjeta ``Estado'' textual & M2, M6, M8 & Datos ya en plot \\
Recordatorio TP/DP texto & M3 & Menú + badge \\
\bottomrule
\end{tabular}
\end{center}
\vspace{0.4cm}

\begin{warnBox}[Caso límite a preservar]
M1 (Coordenadas, $N_i$, Jacobiano) tiene un selector Q4/Q9 \emph{justificado}:
el módulo enseña la diferencia entre interpolación bilineal y bicuadrática.
Es la única excepción legítima. NO muta el proyecto, solo cambia la
visualización.
\end{warnBox}

% ============================================================
\section{Arquitectura propuesta: tres modos de presentación}

La idea central es romper el monopolio del Toplevel y abrir tres canales
para que cada módulo elija el más adecuado.

\subsection{Modo A - Toplevel autónomo (el actual)}

Ventana independiente, ideal cuando el módulo necesita:
\begin{itemize}[leftmargin=1.2em, itemsep=2pt]
  \item Espacio amplio (M4 con notebook 2 sub-tabs).
  \item Comparación lado-a-lado independiente del modelo (M9 sandbox).
  \item Vista 3D de pantalla completa (M7).
\end{itemize}

\subsection{Modo B - Overlay sobre el canvas compartido}

\textbf{Modo nuevo.} El módulo NO abre ventana: se dibuja \emph{encima} del
\texttt{MeshCanvas} principal como un panel translúcido o sobre el propio
canvas (anotaciones, halos, flechas). El alumno sigue interactuando con la
malla normal pero ve la capa educativa superpuesta.

\begin{center}
\begin{tikzpicture}[scale=0.85]
  % main window mockup
  \draw[win] (0,0) rectangle (15, 7);
  \fill[panelBg, rounded corners=2pt] (0.05, 6.4) rectangle (14.95, 6.9);
  \node[text=white, font=\small\bfseries, anchor=west] at (0.3, 6.65) {EduFEM};

  % left tabs
  \draw[panel] (0.2, 0.3) rectangle (4.8, 6.3);
  \fill[phaseProc, rounded corners=1pt] (0.4, 5.85) rectangle (4.6, 6.18);
  \node[text=white, font=\small] at (2.5, 6.0) {PROCESO};

  % educational accordion in the side panel
  \draw[panel, fill=accent!8] (0.4, 4.4) rectangle (4.6, 5.7);
  \node[text=accent, font=\small\bfseries, anchor=north west]
    at (0.55, 5.6) {M2 \textbar{} Matriz B};
  \node[text=white, font=\tiny, anchor=north west]
    at (0.55, 5.25) {Click en el canvas};
  \node[text=white, font=\tiny, anchor=north west]
    at (0.55, 5.0) {para evaluar B en};
  \node[text=white, font=\tiny, anchor=north west]
    at (0.55, 4.75) {un punto Gauss};
  \draw[btnPhase] (0.55, 4.5) rectangle (2.4, 4.7);
  \node[text=white, font=\tiny] at (1.5, 4.6) {Cerrar overlay};

  % canvas with overlay
  \draw[panel] (5.0, 0.3) rectangle (14.8, 6.3);
  % mesh
  \begin{scope}[shift={(6.0, 1.0)}, scale=0.7]
    \foreach \x in {0,1,2,3,4} {
      \foreach \y in {0,1,2,3} {
        \draw[softGray, line width=0.3pt]
          (\x,\y) rectangle ({\x+1},{\y+1});
      }
    }
    \foreach \x in {0,...,5} {
      \foreach \y in {0,...,4} {
        \node[nodeDot] at (\x,\y) {};
      }
    }
    % highlighted element with gauss points
    \fill[selYellow, opacity=0.20] (2,1) rectangle (3,2);
    \draw[selYellow, line width=1pt] (2,1) rectangle (3,2);
    \foreach \pt in {(2.21,1.21), (2.79,1.21), (2.21,1.79), (2.79,1.79)} {
      \node[draw=danger, line width=0.5pt, fill=danger,
            cross out, minimum size=4pt] at \pt {};
    }
  \end{scope}

  % overlay card top right
  \draw[panel, fill=bgDark!90, line width=0.6pt, rounded corners=2pt,
        drop shadow={shadow xshift=1pt, shadow yshift=-1pt}]
    (10.4, 4.7) rectangle (14.5, 5.95);
  \node[text=accent, font=\tiny\bfseries, anchor=north west]
    at (10.55, 5.85) {Matriz B (3x8) en pg1};
  \node[text=white, font=\tiny\ttfamily, anchor=north west]
    at (10.55, 5.55) {[+0.31  0.00  -0.31  0.00 ...]};
  \node[text=white, font=\tiny\ttfamily, anchor=north west]
    at (10.55, 5.30) {[ 0.00 +0.42   0.00 -0.42 ...]};
  \node[text=white, font=\tiny\ttfamily, anchor=north west]
    at (10.55, 5.05) {[+0.42 +0.31  -0.42 -0.31 ...]};
  \node[text=textMuted, font=\tiny, anchor=north west]
    at (10.55, 4.80) {(arrastrable)};

  \node[text=textMuted, font=\tiny, anchor=south] at (7.5, -0.4)
    {Fig. 2 - Modo B: overlay flotante sobre el MeshCanvas compartido};
\end{tikzpicture}
\end{center}
\vspace{0.6cm}

\textbf{Ventajas pedagógicas:}
\begin{itemize}[leftmargin=1.2em, itemsep=2pt]
  \item El alumno NO pierde el contexto del modelo completo.
  \item La selección del canvas (elemento, nodo, arista) es la entrada natural.
  \item La iluminación / coloreado se aplica al canvas real, no a una copia.
\end{itemize}

\subsection{Modo C - Embed en el panel lateral de la fase}

Para módulos cortos y secundarios. El contenido del módulo aparece dentro
de la sub-pestaña ``Educación'' del panel izquierdo (\texttt{pre\_tab},
\texttt{proc\_tab}, \texttt{post\_tab}), sin ventana nueva. El canvas
principal sigue mostrando la malla y reacciona a las acciones del módulo.

\subsection{Tabla de asignación módulo $\rightarrow$ modo}

\begin{center}
\renewcommand{\arraystretch}{1.25}
\begin{tabular}{@{}c l c l@{}}
\toprule
\textbf{Mód.} & \textbf{Concepto} & \textbf{Modo propuesto} & \textbf{Justificación} \\
\midrule
M0 & Calidad de malla        & B - Overlay   & Mapa colorea el canvas real \\
M1 & Coords, $N_i$, Jacobiano & A - Toplevel  & 4 paneles + comparación Q4/Q9 \\
M2 & Matriz B                & B - Overlay   & B asociada a un punto del canvas \\
M3 & Matriz D                & C - Embed     & Cambia con el material elegido \\
M4 & Rigidez + Gauss         & A - Toplevel  & Notebook + cinta de PG \\
M5 & Ensamblaje              & A - Toplevel  & Vuelo elem $\rightarrow$ K global \\
M6 & Fuerzas equivalentes    & B - Overlay   & Acción sobre arista del canvas \\
M7 & Discontinuidad $\sigma$ & A - Toplevel  & Vista 3D en pantalla completa \\
M8 & Direcciones $\sigma_1/\sigma_2$ & B - Overlay & Cruces sobre el canvas + Mohr \\
M9 & Q4 vs Q9 sandbox        & A - Toplevel  & Sandbox independiente \\
\bottomrule
\end{tabular}
\end{center}

\subsection{Operación desde la fase: el panel lateral como controlador}

Para los módulos en modo B (overlay) y C (embed), el panel lateral de la
fase activa actúa como \emph{cabina de comando}. Esto rompe la dependencia
del sidebar interno del Toplevel.

\begin{center}
\begin{tikzpicture}[scale=0.9]
  % phase tab control card
  \draw[panel, fill=panelBg!40] (0, 0) rectangle (6, 6.5);
  \fill[phaseProc, rounded corners=1pt] (0.1, 6.0) rectangle (5.9, 6.4);
  \node[text=white, font=\small\bfseries] at (3.0, 6.2) {PROCESO};

  % educational accordion
  \draw[softGray, line width=0.3pt] (0.2, 5.85) -- (5.8, 5.85);
  \node[text=accent, font=\small\bfseries, anchor=north west]
    at (0.3, 5.75) {Educacion};

  % module list
  \foreach \i/\name/\state in {0/M1 - Coordenadas/normal,
                               1/M2 - Matriz B/active,
                               2/M3 - Matriz D/normal,
                               3/M4 - Rigidez/normal,
                               4/M5 - Ensamblaje/normal,
                               5/M6 - Fuerzas eq./normal} {
    \pgfmathsetmacro{\yy}{5.4 - 0.55*\i}
    \ifthenelse{\equal{\state}{active}}{
      \draw[panel, fill=accent!22] (0.3, \yy) rectangle (5.7, {\yy + 0.45});
      \node[text=white, font=\small\bfseries, anchor=west]
        at (0.45, {\yy + 0.22}) {\name};
      \node[text=accent, font=\tiny, anchor=east]
        at (5.55, {\yy + 0.22}) {ACTIVO};
    }{
      \draw[softGray, line width=0.2pt] (0.3, \yy) rectangle (5.7, {\yy + 0.45});
      \node[text=textMuted, font=\small, anchor=west]
        at (0.45, {\yy + 0.22}) {\name};
    }
  }

  % active module dynamic controls
  \draw[softGray, line width=0.3pt] (0.2, 2.0) -- (5.8, 2.0);
  \node[text=accent, font=\small\bfseries, anchor=north west]
    at (0.3, 1.9) {Controles M2};
  \node[text=textMuted, font=\tiny, anchor=north west]
    at (0.3, 1.55) {Punto de evaluacion};
  \draw[btn] (0.3, 1.0) rectangle (1.2, 1.3);  \node[text=white, font=\tiny] at (0.75, 1.15) {centro};
  \draw[btn] (1.3, 1.0) rectangle (2.2, 1.3);  \node[text=white, font=\tiny] at (1.75, 1.15) {pg1};
  \draw[btn] (2.3, 1.0) rectangle (3.2, 1.3);  \node[text=white, font=\tiny] at (2.75, 1.15) {pg2};
  \draw[btn] (3.3, 1.0) rectangle (4.2, 1.3);  \node[text=white, font=\tiny] at (3.75, 1.15) {pg3};
  \draw[btn] (4.3, 1.0) rectangle (5.2, 1.3);  \node[text=white, font=\tiny] at (4.75, 1.15) {pg4};

  \draw[btnPhase] (0.3, 0.4) rectangle (3.0, 0.7);
  \node[text=white, font=\tiny] at (1.65, 0.55) {Iluminar elem. en canvas};
  \draw[btn] (3.1, 0.4) rectangle (5.7, 0.7);
  \node[text=white, font=\tiny] at (4.4, 0.55) {Mostrar B en LaTeX};

  \node[text=textMuted, font=\tiny, anchor=south] at (3, -0.4)
    {Fig. 3 - Panel lateral de la fase como cabina de comando};
\end{tikzpicture}
\end{center}
\vspace{0.4cm}

\textbf{Reglas de oro de este patrón:}
\begin{enumerate}[leftmargin=1.2em, itemsep=2pt]
  \item El \textbf{elemento bajo análisis} se elige por \emph{click en el
        MeshCanvas}, no por combobox interno.
  \item La \textbf{iluminación} (highlight, glow, halo) se aplica al
        elemento real del canvas. El módulo solo solicita ``ilumina el
        elemento $E_3$ en color X'' y el canvas obedece.
  \item Los \textbf{controles dinámicos} del módulo aparecen en el panel
        lateral, sin abrir ventana. Para datos extensos (matrices,
        gráficos secundarios), un \emph{overlay flotante} sobre el
        canvas — arrastrable y cerrable.
\end{enumerate}

\newpage

% ============================================================
\section{Propuestas por módulo}

A continuación se desarrolla cada módulo con: concepto enseñado, problema
del layout actual, mockup del rediseño y justificación pedagógica.

% ============================================================
\subsection{M0 \textbar{} Calidad geométrica de la malla}

\begin{moduleBox}[M0 - Calidad geometrica - Modo B (Overlay)]{phasePre}

\textbf{Concepto enseñado.} Métricas de calidad de elementos cuadriláteros
(Jacobian Ratio, Robinson Aspect, Taper / Admissibility, Scaled Jacobian)
y cómo degradan al distorsionar un nodo.

\textbf{Patología actual.}
\begin{itemize}[leftmargin=1em]
  \item Los criterios de calidad están en texto en el sidebar (lectura única).
  \item La ``tarjeta de métricas'' duplica el radar.
  \item Click sobre un elemento en el panel ``mapa'' es la única interacción.
\end{itemize}

\textbf{Propuesta UX - Vista rayos X sobre el canvas:}
\begin{enumerate}[leftmargin=1.2em, itemsep=2pt]
  \item El \texttt{MeshCanvas} principal pinta cada elemento con verde /
        amarillo / rojo según su estado combinado. Sin abrir ventana.
  \item \textbf{Hover sobre un elemento}: aparece una tarjeta flotante
        anclada al cursor con el radar de 4 ejes.
  \item \textbf{Drag sobre un nodo}: distorsión interactiva en tiempo real,
        métricas recalculan a 60\,fps.
  \item Histograma horizontal colapsable de 50\,px en el footer del canvas
        (no en el sidebar).
\end{enumerate}

\begin{center}
\begin{tikzpicture}[scale=0.92]
  \draw[win] (0,0) rectangle (15, 6.5);

  % canvas grid colored by quality
  \begin{scope}[shift={(0.5, 0.5)}]
    \foreach \r/\c/\col in {%
      0/0/good, 0/1/good, 0/2/warn,   0/3/good,
      1/0/good, 1/1/warn, 1/2/danger, 1/3/warn,
      2/0/good, 2/1/good, 2/2/good,   2/3/good} {
      \fill[\col, opacity=0.5]
        ({\c*1.4}, {\r*1.4}) rectangle ({\c*1.4 + 1.4}, {\r*1.4 + 1.4});
      \draw[white, line width=0.3pt]
        ({\c*1.4}, {\r*1.4}) rectangle ({\c*1.4 + 1.4}, {\r*1.4 + 1.4});
      \node[nodeDot] at ({\c*1.4}, {\r*1.4}) {};
      \node[nodeDot] at ({\c*1.4 + 1.4}, {\r*1.4}) {};
    }
    % top row of nodes
    \foreach \c in {0,1,2,3,4} {
      \node[nodeDot] at ({\c*1.4}, {3*1.4}) {};
    }
    % distorted node (drag preview)
    \node[selRing] at (3.7, 2.55) {};
    \draw[selYellow, dashed, line width=0.4pt, ->]
      (2.8, 2.8) -- (3.7, 2.55);
  \end{scope}

  % floating radar card
  \draw[panel, fill=bgDark!95, line width=0.6pt, rounded corners=2pt,
        drop shadow={shadow xshift=1pt, shadow yshift=-1pt}]
    (8.5, 1.5) rectangle (14.5, 5.5);
  \node[text=danger, font=\small\bfseries, anchor=north west]
    at (8.65, 5.4) {Elemento E6 - Mala};

  % radar diagram
  \begin{scope}[shift={(11.5, 3.3)}, scale=0.85]
    \foreach \r in {0.4, 0.8, 1.2} {
      \draw[softGray, line width=0.3pt] (0,0) circle (\r);
    }
    \foreach \a in {0, 90, 180, 270} {
      \draw[softGray, line width=0.3pt] (0,0) -- (\a:1.3);
    }
    \node[text=white, font=\tiny] at (0:1.5)   {SJ};
    \node[text=white, font=\tiny] at (90:1.5)  {JR};
    \node[text=white, font=\tiny] at (180:1.5) {AR};
    \node[text=white, font=\tiny] at (270:1.5) {Tap};
    \fill[danger, opacity=0.4]
      (0:0.5) -- (90:0.7) -- (180:0.4) -- (270:0.3) -- cycle;
    \draw[danger, line width=0.6pt]
      (0:0.5) -- (90:0.7) -- (180:0.4) -- (270:0.3) -- cycle;
  \end{scope}

  \node[text=white, font=\tiny\ttfamily, anchor=north west]
    at (8.65, 2.0) {SJ=0.18  JR=0.28  AR=4.2};

  % footer histogram strip
  \draw[panel, fill=panelBg!50] (0.5, -0.05) rectangle (14.5, 0.4);
  \node[text=accent, font=\tiny, anchor=west] at (0.6, 0.18) {Distrib. Jacobian};
  \foreach \x/\h in {2.5/0.08, 3.0/0.16, 3.5/0.22, 4.0/0.30, 4.5/0.30,
                     5.0/0.18, 5.5/0.10, 6.0/0.06} {
    \fill[accent, opacity=0.6] (\x, 0.06) rectangle ({\x+0.4}, {0.06 + \h});
  }

  \node[text=textMuted, font=\tiny, anchor=south] at (7.5, -0.55)
    {Fig. 4 - M0: hover sobre elemento abre radar flotante; drag de nodo distorsiona en vivo};
\end{tikzpicture}
\end{center}

\textbf{Justificación pedagógica.} El alumno ve la calidad ``en su sitio'' —
sobre la malla real que está construyendo o cargando, no sobre una copia en
panel auxiliar. El drag interactivo convierte la métrica en algo
\emph{cinético}: el alumno ve el Jacobian caer al mover el nodo, y eso es
el aprendizaje real, no una tabla de números.

\end{moduleBox}

% ============================================================
\subsection{M1 \textbar{} Coordenadas naturales, $N_i$, Jacobiano}

\begin{moduleBox}[M1 - Mapeo isoparametrico - Modo A (Toplevel)]{phaseProc}

\textbf{Concepto enseñado.} El mapeo isoparamétrico
$(x, y) \leftrightarrow (\xi, \eta)$, las funciones de forma $N_i$ y el
Jacobiano que conecta ambos espacios.

\textbf{Patología actual.} Cuatro paneles 2$\times$2 + sidebar denso. La
relación entre el panel físico y el panel natural se pierde porque el
alumno tiene que mirar dos esquinas opuestas.

\textbf{Propuesta UX - Split mirror animado:}
\begin{enumerate}[leftmargin=1.2em, itemsep=2pt]
  \item Dos mitades simétricas separadas por un eje vertical con la
        flecha ``$N_i \rightarrow$''. Físico izq, natural der.
  \item \textbf{Hover sin click} en cualquier mitad dibuja un punto rojo
        sincronizado en ambas y también en las superficies $N_i$ y
        det\,$\mathbf{J}$ (insets pequeños, esquinas).
  \item Chips numerados 1..9 sobre los nodos para elegir $N_i$.
  \item \textbf{Demo-mode al abrir}: una ``mosca'' recorre las 9 funciones
        de forma en 5\,s.
\end{enumerate}

\begin{center}
\begin{tikzpicture}[scale=0.85]
  \draw[win] (0,0) rectangle (16, 7.5);

  % header
  \fill[panelBg, rounded corners=2pt] (0.05, 6.95) rectangle (15.95, 7.45);
  \node[text=accent, font=\small\bfseries, anchor=west]
    at (0.3, 7.2) {M1 - Mapeo isoparametrico};
  % chips for N_i
  \foreach \i in {1,...,9} {
    \draw[chip] ({3.5 + 0.5*\i}, 7.05) rectangle ({3.85 + 0.5*\i}, 7.35);
    \node[text=white, font=\tiny] at ({3.675 + 0.5*\i}, 7.20) {\i};
  }

  % left half: physical
  \draw[panel] (0.2, 0.3) rectangle (7.7, 6.7);
  \node[text=textMuted, font=\tiny, anchor=south west]
    at (0.4, 6.45) {Espacio fisico (x, y)};
  \begin{scope}[shift={(2.0, 1.5)}, scale=1.1]
    \draw[accent, fill=accent!10, line width=1pt]
      (0,0) -- (3,0.2) -- (3.2,2.5) -- (0.1,2.3) -- cycle;
    \node[nodeDot] at (0,0) {};         \node[label] at (-0.2,-0.2) {1};
    \node[nodeDot] at (3,0.2) {};       \node[label] at (3.2,0.0) {2};
    \node[nodeDot] at (3.2,2.5) {};     \node[label] at (3.4,2.6) {3};
    \node[nodeDot] at (0.1,2.3) {};     \node[label] at (-0.2,2.4) {4};
    \node[selRing] at (1.6, 1.2) {};
  \end{scope}

  % central arrow
  \draw[accent, line width=1pt, ->] (7.85, 3.5) -- (8.35, 3.5);
  \node[text=accent, font=\tiny] at (8.1, 3.85) {$N_i(\xi,\eta)$};

  % right half: natural
  \draw[panel] (8.5, 0.3) rectangle (16, 6.7);
  \node[text=textMuted, font=\tiny, anchor=south west]
    at (8.7, 6.45) {Espacio natural (xi, eta)};
  \begin{scope}[shift={(11.0, 2.0)}, scale=1.5]
    \draw[accent, fill=accent!10, line width=1pt]
      (-1, -1) rectangle (1, 1);
    \draw[softGray, line width=0.3pt] (-1, 0) -- (1, 0);
    \draw[softGray, line width=0.3pt] (0, -1) -- (0, 1);
    \node[nodeDot] at (-1,-1) {}; \node[label] at (-1.2,-1.1) {1};
    \node[nodeDot] at (1,-1) {};  \node[label] at (1.2,-1.1) {2};
    \node[nodeDot] at (1,1) {};   \node[label] at (1.2,1.1) {3};
    \node[nodeDot] at (-1,1) {};  \node[label] at (-1.2,1.1) {4};
    \node[selRing] at (0.1, 0.05) {};
  \end{scope}

  % inset: N_i surface
  \draw[panel, fill=bgDark!90, rounded corners=2pt, line width=0.4pt]
    (10.2, 4.7) rectangle (12.4, 6.4);
  \node[text=textMuted, font=\tiny, anchor=north west]
    at (10.3, 6.35) {N1(xi, eta)};
  \draw[accent, line width=0.4pt] plot[smooth, samples=20, domain=-1:1]
    ({11.3 + 0.4*\x}, {5.2 + 0.7*(1-\x)*(1-\x)*0.25});

  % inset: detJ
  \draw[panel, fill=bgDark!90, rounded corners=2pt, line width=0.4pt]
    (13.0, 4.7) rectangle (15.5, 6.4);
  \node[text=textMuted, font=\tiny, anchor=north west]
    at (13.1, 6.35) {det J(xi, eta)};
  \draw[warn, line width=0.4pt] plot[smooth, samples=20, domain=-1:1]
    ({14.25 + 0.5*\x}, {5.5 + 0.3*\x});

  \node[text=textMuted, font=\tiny, anchor=south] at (8, -0.4)
    {Fig. 5 - M1: split mirror con marcador rojo sincronizado en ambos espacios};
\end{tikzpicture}
\end{center}

\textbf{Justificación pedagógica.} El isomorfismo $\xi,\eta \leftrightarrow x,y$
se aprende mejor cuando ambos lados se mueven \emph{a la vez}. El mouse
es un puente físico entre los dos mundos. Los insets $N_i$ y det\,$\mathbf{J}$
ya no compiten por espacio: son referencias contextuales.

\end{moduleBox}

% ============================================================
\subsection{M2 \textbar{} Matriz B - Modo B (Overlay)}

\begin{moduleBox}[M2 - Matriz B - Modo B (Overlay)]{phaseProc}

\textbf{Concepto enseñado.} La matriz $\mathbf{B}$ que relaciona
desplazamientos nodales con deformaciones $\boldsymbol{\varepsilon} =
\mathbf{B}\mathbf{u}$, su construcción vía
$\partial N_i / \partial x = \mathbf{J}^{-1} \cdot \partial N_i / \partial \xi$
y la propiedad de superconvergencia en puntos de Gauss.

\textbf{Patología actual.} La matriz $\mathbf{B}$ aparece en LaTeX abajo,
pero no se conecta visualmente con cuál término viene de cuál nodo. La
``superconvergencia'' es una palabra del sidebar.

\textbf{Propuesta UX - Matriz como protagonista, geometría secundaria:}
\begin{enumerate}[leftmargin=1.2em, itemsep=2pt]
  \item El módulo NO abre Toplevel. Se activa desde la sub-pestaña
        ``Educación'' del panel \texttt{proc\_tab}.
  \item En el canvas: glow amarillo pulsante en los puntos de Gauss del
        elemento seleccionado.
  \item \textbf{Click en un punto Gauss} sobre el canvas $\rightarrow$
        aparece un \emph{overlay flotante} en el canvas con la matriz
        $\mathbf{B}$ en LaTeX grande, evaluada en ese punto.
  \item \textbf{Hover sobre una columna} de $\mathbf{B}$ ilumina el nodo
        correspondiente en la malla; \textbf{hover sobre una fila} ilumina
        la componente de deformación ($\varepsilon_x$, $\varepsilon_y$,
        $\gamma_{xy}$).
\end{enumerate}

\begin{center}
\begin{tikzpicture}[scale=0.85]
  \draw[win] (0,0) rectangle (16, 7.5);

  % left: phase tab
  \draw[panel] (0.2, 0.3) rectangle (4.8, 7.2);
  \fill[phaseProc, rounded corners=1pt] (0.4, 6.7) rectangle (4.6, 7.05);
  \node[text=white, font=\small\bfseries] at (2.5, 6.875) {PROCESO};
  \node[text=accent, font=\tiny\bfseries, anchor=north west]
    at (0.4, 6.5) {Educacion};
  \draw[panel, fill=accent!22] (0.4, 5.5) rectangle (4.6, 6.2);
  \node[text=white, font=\small, anchor=west] at (0.5, 5.85)
    {M2 - Matriz B (activo)};

  \node[text=textMuted, font=\tiny, anchor=west] at (0.5, 5.25)
    {Click en pg amarillos};

  % canvas
  \draw[panel] (5.0, 0.3) rectangle (15.8, 7.2);
  \begin{scope}[shift={(6.0, 1.0)}, scale=0.6]
    \foreach \x in {0,1,2,3,4} {
      \foreach \y in {0,1,2,3} {
        \draw[softGray, line width=0.3pt]
          (\x,\y) rectangle ({\x+1},{\y+1});
      }
    }
    \foreach \x in {0,...,5} {
      \foreach \y in {0,...,4} {
        \node[nodeDot] at (\x,\y) {};
      }
    }
    \fill[selYellow, opacity=0.18] (2,1) rectangle (3,2);
    \draw[selYellow, line width=1pt] (2,1) rectangle (3,2);
    % Glow pulsante en pg
    \foreach \pt/\sel in {(2.21,1.21)/0, (2.79,1.21)/1,
                          (2.21,1.79)/0, (2.79,1.79)/0} {
      \node[circle, draw=selYellow, line width=0.6pt,
            fill=selYellow!40, minimum size=8pt, inner sep=0pt] at \pt {};
      \ifnum\sel=1
        \node[circle, draw=selYellow, line width=1pt,
              minimum size=14pt, inner sep=0pt] at \pt {};
      \fi
    }
  \end{scope}

  % overlay matrix card
  \draw[panel, fill=bgDark!95, rounded corners=2pt, line width=0.7pt,
        drop shadow={shadow xshift=1pt, shadow yshift=-1pt}]
    (8.5, 3.2) rectangle (15.5, 6.7);
  \node[text=accent, font=\small\bfseries, anchor=north west]
    at (8.65, 6.55) {Matriz B (3x8) en pg2};
  \node[text=warn, font=\tiny, anchor=north west]
    at (8.65, 6.25) {SUPERCONVERGENCIA: sigma exacta aqui};

  % matrix LaTeX-like
  \node[text=white, font=\tiny\ttfamily, anchor=north west]
    at (8.65, 5.85)
    {\begin{tabular}{cccccccc}
      $+$0.31 & 0 & $-$0.31 & 0 & $+$0.31 & 0 & $-$0.31 & 0 \\
      0 & $-$0.42 & 0 & $-$0.42 & 0 & $+$0.42 & 0 & $+$0.42 \\
      $-$0.42 & $+$0.31 & $-$0.42 & $-$0.31 & $+$0.42 & $+$0.31 & $+$0.42 & $-$0.31 \\
    \end{tabular}};

  % rows labels
  \node[text=danger, font=\tiny, anchor=east] at (8.55, 5.65) {eps\_x};
  \node[text=accent, font=\tiny, anchor=east] at (8.55, 5.50) {eps\_y};
  \node[text=warn,   font=\tiny, anchor=east] at (8.55, 5.30) {gamma\_xy};

  % column labels (nodes)
  \foreach \i/\xx in {1/9.05, 2/9.95, 3/10.85, 4/11.75, 5/12.65, 6/13.55, 7/14.45, 8/15.30} {
    \node[text=textMuted, font=\tiny] at (\xx, 5.05) {N\i};
  }

  \node[text=textMuted, font=\tiny, anchor=north west]
    at (8.65, 4.30) {Hover en columna -> ilumina nodo};
  \node[text=textMuted, font=\tiny, anchor=north west]
    at (8.65, 4.05) {Hover en fila -> ilumina componente};

  \node[text=textMuted, font=\tiny, anchor=south] at (8, -0.4)
    {Fig. 6 - M2: overlay con matriz B grande, vinculacion bidireccional con la malla};
\end{tikzpicture}
\end{center}

\textbf{Justificación pedagógica.} La superconvergencia deja de ser un
término abstracto: el alumno \emph{ve} que esos puntos amarillos pulsantes
son los que el solver elige y que la $\mathbf{B}$ que vive ahí es la que
importa para $\sigma$. Las flechas hover-driven entre columnas y nodos
materializan la regla \emph{cada par de columnas pertenece a un nodo}.

\end{moduleBox}

% ============================================================
\subsection{M3 \textbar{} Matriz constitutiva $\mathbf{D}$ - Modo C (Embed)}

\begin{moduleBox}[M3 - Matriz D - Modo C (Embed en panel lateral)]{phaseProc}

\textbf{Concepto enseñado.} La matriz $\mathbf{D}(E, \nu, \text{caso})$
que conecta tensión y deformación $\boldsymbol{\sigma} = \mathbf{D}
\boldsymbol{\varepsilon}$, dependiendo del material y del caso plano.

\textbf{Patología actual.} La ventana es enorme para algo que solo muestra
una matriz $3 \times 3$ y una fórmula. El selector de elemento se duplica.
El ``recordatorio teórico'' es texto estático.

\textbf{Propuesta UX - Sin Toplevel: panel embebido + dial físico:}
\begin{enumerate}[leftmargin=1.2em, itemsep=2pt]
  \item M3 vive dentro del panel \texttt{proc\_tab\,$\triangleright$\,Educación}.
        El elemento se elige con click sobre el canvas (highlight amarillo).
  \item \textbf{Tres tarjetas verticales} en el panel: $\mathbf{D}$ simbólica,
        $\mathbf{D}$ numérica, animación esquemática TP vs DP.
  \item Único control: un \textbf{dial circular de Poisson} ($\nu \in
        [0, 0.499]$). Mover el dial reescribe $\mathbf{D}$ en tiempo real.
        El dial se vuelve rojo cuando $\nu \to 0.5$ en DP (volumetric locking).
  \item \textbf{Animación TP vs DP}: una placa fina que se carga lateralmente
        (TP) vs un muro de contención infinito (DP), con el cubito 2D
        deformándose de modo distinto.
\end{enumerate}

\begin{center}
\begin{tikzpicture}[scale=0.92]
  % lateral panel
  \draw[panel, fill=panelBg!30] (0, 0) rectangle (5.2, 7.5);
  \fill[phaseProc, rounded corners=1pt] (0.1, 7.05) rectangle (5.1, 7.40);
  \node[text=white, font=\small\bfseries] at (2.6, 7.225) {PROCESO};
  \node[text=accent, font=\small\bfseries, anchor=north west]
    at (0.2, 6.95) {M3 - Matriz D};

  % dial
  \begin{scope}[shift={(2.6, 5.5)}]
    \draw[softGray, line width=2pt] (0,0) circle (0.85);
    \draw[accent, line width=2.5pt]
      (0.85, 0) arc (0:140:0.85);
    \node[circle, fill=warn, draw=white, minimum size=8pt, inner sep=0pt]
      at ({140}:0.85) {};
    \node[text=white, font=\small\bfseries] at (0,0) {$\nu$};
    \node[text=warn, font=\small\bfseries] at (0,-0.45) {0.30};
  \end{scope}

  % D symbolic card
  \draw[panel, fill=bgDark!90] (0.2, 3.4) rectangle (5.1, 4.5);
  \node[text=textMuted, font=\tiny, anchor=north west]
    at (0.3, 4.4) {D simbolica (TP)};
  \node[text=white, font=\small\ttfamily, anchor=north west]
    at (0.3, 4.05) {E/(1-v\textasciicircum 2) [1 v 0; v 1 0; 0 0 (1-v)/2]};

  % D numeric card
  \draw[panel, fill=bgDark!90] (0.2, 1.7) rectangle (5.1, 3.2);
  \node[text=textMuted, font=\tiny, anchor=north west]
    at (0.3, 3.1) {D numerica (Pa)};
  \node[text=accent, font=\tiny\ttfamily, anchor=north west]
    at (0.3, 2.85) {2.31e11  6.92e10  0};
  \node[text=accent, font=\tiny\ttfamily, anchor=north west]
    at (0.3, 2.55) {6.92e10  2.31e11  0};
  \node[text=accent, font=\tiny\ttfamily, anchor=north west]
    at (0.3, 2.25) {     0       0   8.08e10};

  % TP/DP animation card
  \draw[panel, fill=bgDark!90] (0.2, 0.2) rectangle (5.1, 1.5);
  \node[text=textMuted, font=\tiny, anchor=north west]
    at (0.3, 1.4) {Caso plano: TP (placa delgada)};
  \fill[accent!40] (0.6, 0.5) rectangle (2.0, 1.0);
  \draw[white, line width=0.4pt] (0.6, 0.5) rectangle (2.0, 1.0);
  \draw[danger, line width=0.6pt, ->] (2.2, 0.75) -- (2.6, 0.75);
  \draw[danger, line width=0.6pt, ->] (0.4, 0.75) -- (0.0, 0.75);
  \node[text=accent, font=\tiny] at (3.5, 0.75) {sigma\_z = 0};

  % canvas right
  \draw[panel] (5.5, 0.3) rectangle (16, 7.5);
  \node[text=textMuted, font=\tiny, anchor=north west]
    at (5.7, 7.4) {MeshCanvas (elemento clickeado se ilumina aqui)};
  \begin{scope}[shift={(7, 1.5)}, scale=0.7]
    \foreach \x in {0,1,2,3,4} {
      \foreach \y in {0,1,2,3} {
        \draw[softGray, line width=0.3pt]
          (\x,\y) rectangle ({\x+1},{\y+1});
        \node[nodeDot] at (\x,\y) {};
      }
    }
    \foreach \x in {0,1,2,3,4,5} {
      \node[nodeDot] at (\x,4) {};
    }
    \foreach \y in {0,...,4} {
      \node[nodeDot] at (5,\y) {};
    }
    \fill[selYellow, opacity=0.30] (2,1) rectangle (3,2);
    \draw[selYellow, line width=1.5pt] (2,1) rectangle (3,2);
    \node[text=selYellow, font=\small\bfseries] at (2.5, 1.5) {E6};
  \end{scope}
  \draw[selYellow, dashed, line width=0.4pt, ->]
    (4.9, 5.5) -- (8.95, 2.55);
  \node[text=selYellow, font=\tiny] at (6.8, 4.3) {material -> D};

  \node[text=textMuted, font=\tiny, anchor=south] at (8, -0.4)
    {Fig. 7 - M3: panel embebido con dial de Poisson; click en canvas elige el elemento};
\end{tikzpicture}
\end{center}

\textbf{Justificación pedagógica.} El dial es \emph{cinético}: el alumno
experimenta la singularidad en $\nu = 0.5$ moviendo la perilla y viendo
los números explotar. El cubito animado da intuición física (TP = placa,
DP = muro). M3 es ahora un módulo \emph{secundario} acoplado al canvas,
no un Toplevel intrusivo.

\end{moduleBox}

% ============================================================
\subsection{M4 \textbar{} Rigidez del elemento + Cuadratura de Gauss}

\begin{moduleBox}[M4 - Rigidez + Gauss - Modo A (Toplevel)]{phaseProc}

\textbf{Concepto enseñado.} $\mathbf{k}_e = \int \int \mathbf{B}^T
\mathbf{D} \mathbf{B} |\det \mathbf{J}| t \, d\xi \, d\eta$ aproximado
por cuadratura de Gauss. El integrando simbólico es feroz; la
cuadratura es elegante.

\textbf{Patología actual.} Sidebar con 3 spinboxes ($E$, $\nu$, $t$) que
\emph{no deberían existir} (vienen del material). Notebook con dos
sub-tabs separadas, sin transición narrativa entre ellas.

\textbf{Propuesta UX - Cinta transportadora ``fábrica de K'':}
\begin{enumerate}[leftmargin=1.2em, itemsep=2pt]
  \item Eliminar spinboxes $E$, $\nu$, $t$. Se leen del material.
  \item Reemplazar las dos sub-tabs por una \textbf{cinta horizontal} en
        la mitad inferior: cada PG entra desde la izquierda con su
        contribución $w_p \mathbf{B}^T\mathbf{D}\mathbf{B} |J| t$ y se
        suma al heatmap acumulado de $\mathbf{k}_e$ a la derecha.
  \item \textbf{Scrubber temporal} debajo (timeline). Arrastrar rebobina
        el ensamblado.
  \item \textbf{Único control}: orden de cuadratura $1\times1, 2\times2,
        3\times3$. Al elegir $1\times1$ aparecen los \textbf{modos hourglass}
        sobre la malla — momento didáctico clave (sub-integración).
\end{enumerate}

\begin{center}
\begin{tikzpicture}[scale=0.9]
  \draw[win] (0,0) rectangle (16, 7);

  % top half: integrando simbolico + superficie K_ij
  \draw[panel] (0.3, 3.5) rectangle (8, 6.8);
  \node[text=textMuted, font=\tiny, anchor=north west]
    at (0.4, 6.7) {Integrando K\_ij(xi, eta) en LaTeX};
  \node[text=white, font=\footnotesize\ttfamily, anchor=north west]
    at (0.4, 6.3)
    {K\_22 = (E*t/((1-v\^{}2)*detJ\^{}2)) *};
  \node[text=white, font=\footnotesize\ttfamily, anchor=north west]
    at (0.4, 5.95)
    {[(1-eta)\^{}2/16 * J22\^{}2 +};
  \node[text=white, font=\footnotesize\ttfamily, anchor=north west]
    at (0.4, 5.65)
    {(1-xi)\^{}2/16 * J11\^{}2 +};
  \node[text=textMuted, font=\footnotesize, anchor=north west]
    at (0.4, 5.30)
    {... (terminos cruzados) ...};
  \node[text=warn, font=\tiny, anchor=north west]
    at (0.4, 4.85)
    {No hay primitiva cerrada: usar Gauss};

  \draw[panel] (8.2, 3.5) rectangle (15.7, 6.8);
  \node[text=textMuted, font=\tiny, anchor=north west]
    at (8.3, 6.7) {Superficie |K\_ij(xi, eta)| - 3D};
  \begin{scope}[shift={(11.95, 5.0)}, scale=1.0]
    \foreach \i in {0,...,5} {
      \foreach \j in {0,...,5} {
        \pgfmathsetmacro{\zz}{0.6 - 0.04*((\i-2.5)^2 + (\j-2.5)^2)}
        \pgfmathsetmacro{\rr}{0.6*\zz + 0.4}
        \pgfmathsetmacro{\bb}{0.4*\zz}
        \fill[red!30!orange, opacity=\rr]
          ({-1.5 + 0.5*\i}, {-0.6 + 0.3*\j}) rectangle
          ({-1 + 0.5*\i}, {-0.3 + 0.3*\j});
      }
    }
  \end{scope}

  % bottom half: gauss conveyor
  \draw[panel] (0.3, 0.3) rectangle (15.7, 3.3);
  \node[text=accent, font=\small\bfseries, anchor=north west]
    at (0.4, 3.2) {Cinta transportadora 2x2};

  % gauss point pills moving right
  \foreach \i/\xx/\op in {1/2.0/0.4, 2/4.5/0.4, 3/7.0/1.0, 4/9.5/0.25} {
    \draw[selYellow, line width=0.5pt, fill=selYellow, opacity=\op,
          rounded corners=2pt]
      (\xx, 1.6) rectangle ({\xx + 1.6}, 2.6);
    \node[text=white, font=\tiny\bfseries] at ({\xx + 0.8}, 2.30)
      {pg\i};
    \node[text=white, font=\tiny] at ({\xx + 0.8}, 2.0)
      {w*BtDB*detJ*t};
  }

  % arrow conveyor
  \draw[accent, line width=0.6pt, ->] (0.5, 1.2) -- (10.5, 1.2);
  \node[text=accent, font=\tiny] at (5.5, 0.95) {sumar al heatmap K\_e};

  % accumulated K heatmap right
  \draw[panel, fill=bgDark!80] (11.5, 0.5) rectangle (15.5, 3.0);
  \node[text=textMuted, font=\tiny, anchor=north west]
    at (11.6, 2.95) {K\_e acumulado (3 / 4 PG)};
  \begin{scope}[shift={(13.5, 1.7)}]
    \foreach \i in {0,...,4} {
      \foreach \j in {0,...,4} {
        \pgfmathsetmacro{\val}{0.5 + 0.4*sin((\i+\j)*30)}
        \fill[accent, opacity=\val]
          ({-1 + 0.4*\i}, {-1 + 0.4*\j}) rectangle
          ({-0.6 + 0.4*\i}, {-0.6 + 0.4*\j});
      }
    }
  \end{scope}

  \node[text=textMuted, font=\tiny, anchor=south] at (8, -0.4)
    {Fig. 8 - M4: integrando arriba, cinta de PG abajo, K\_e ensamblandose en vivo};
\end{tikzpicture}
\end{center}

\textbf{Justificación pedagógica.} El concepto difícil es ``por qué hay
puntos discretos donde se evalúa''. La metáfora de cinta transportadora
hace que el alumno \emph{cuente} los términos: 4 contribuciones para
$2\times2$, 9 para $3\times3$. La aparición de modos hourglass en
$1\times1$ es un cliffhanger pedagógico que ningún sidebar puede
transmitir.

\end{moduleBox}

% ============================================================
\subsection{M5 \textbar{} Ensamblaje global - Vuelo de elementos}

\begin{moduleBox}[M5 - Ensamblaje - Modo A (Toplevel)]{phaseProc}

\textbf{Concepto enseñado.} Reubicación de cada $\mathbf{k}_e$ local
en las filas y columnas globales correspondientes a sus GDLs, hasta
formar $\mathbf{K}$ globales y aplicar BCs.

\textbf{Patología actual.} La animación ``flying element'' actual hace
solo \emph{fade-in} en las celdas destino. La metáfora del vuelo NO se
materializa: el alumno no ve por qué se llama ``flying''.

\textbf{Propuesta UX - Vuelo real con curvas de Bézier:}
\begin{enumerate}[leftmargin=1.2em, itemsep=2pt]
  \item Al pulsar ``ensamblar siguiente'', el bloque $\mathbf{k}_e$ se
        \emph{despega} de la posición del elemento, se eleva, gira y
        \emph{vuela} en una curva Bézier hasta su bloque destino en
        $\mathbf{K}$.
  \item Trazas Bézier residuales (fade) muestran las migraciones previas.
        Al final el alumno ve un ``mapa de vuelos'' sobre la malla.
  \item \textbf{Click en una fila/columna restringida}: tachado visual
        animado representando la eliminación del GDL.
\end{enumerate}

\begin{center}
\begin{tikzpicture}[scale=0.85]
  \draw[win] (0,0) rectangle (16, 7);

  % left mesh
  \draw[panel] (0.3, 0.3) rectangle (7, 6.7);
  \node[text=textMuted, font=\tiny, anchor=north west]
    at (0.4, 6.6) {Malla - 6 elementos};
  \begin{scope}[shift={(1.0, 0.8)}, scale=0.8]
    \foreach \x in {0,1,2,3} {
      \foreach \y in {0,1} {
        \draw[softGray, line width=0.3pt]
          (\x,\y) rectangle ({\x+1},{\y+1});
        \node[nodeDot] at (\x,\y) {};
      }
    }
    \foreach \x in {0,1,2,3,4} {
      \node[nodeDot] at (\x,2) {};
    }
    \foreach \y in {0,1,2} {
      \node[nodeDot] at (4,\y) {};
    }
    % element about to fly
    \fill[selYellow, opacity=0.5] (2,1) rectangle (3,2);
    \draw[selYellow, line width=1.5pt] (2,1) rectangle (3,2);
    \node[text=white, font=\small\bfseries] at (2.5, 1.5) {E3};
  \end{scope}

  % right K matrix
  \draw[panel] (7.3, 0.3) rectangle (15.7, 6.7);
  \node[text=textMuted, font=\tiny, anchor=north west]
    at (7.4, 6.6) {K global (heatmap)};

  \begin{scope}[shift={(11.5, 3.5)}, scale=0.55]
    \foreach \i in {0,...,9} {
      \foreach \j in {0,...,9} {
        \pgfmathsetmacro{\val}{0.05 + 0.4*((\i==\j) ? 1 : 0) +
                                0.15*((abs(\i-\j) < 3) ? 1 : 0)}
        \fill[accent, opacity=\val]
          ({-3 + 0.6*\i}, {-3 + 0.6*\j}) rectangle
          ({-2.4 + 0.6*\i}, {-2.4 + 0.6*\j});
      }
    }
    % destination cells highlighted (block 4..7)
    \foreach \i in {4,5,6,7} {
      \foreach \j in {4,5,6,7} {
        \draw[selYellow, line width=0.6pt]
          ({-3 + 0.6*\i}, {-3 + 0.6*\j}) rectangle
          ({-2.4 + 0.6*\i}, {-2.4 + 0.6*\j});
      }
    }
    % BC strikethrough on row/col 0
    \draw[danger, line width=1pt]
      (-3, -3.3) -- (-3, 3);
    \draw[danger, line width=1pt]
      (-3.3, -3) -- (3, -3);
  \end{scope}

  % flying ke trajectory (Bezier)
  \draw[selYellow, line width=1pt, ->]
    (3.0, 3.5) .. controls (5, 5.5) and (8, 5) .. (11, 4.0);

  % flying block mid-air
  \draw[selYellow, fill=selYellow!40, line width=0.6pt, rounded corners=1pt,
        rotate around={20:(7.5, 4.7)}]
    (7.0, 4.4) rectangle (8.0, 5.0);
  \node[text=white, font=\tiny\bfseries] at (7.5, 4.7) {k\_e};

  \node[text=textMuted, font=\tiny, anchor=south] at (8, -0.4)
    {Fig. 9 - M5: k\_e literalmente vuela del elemento al bloque destino en K};
\end{tikzpicture}
\end{center}

\textbf{Justificación pedagógica.} El nombre ``flying element'' debe
traducirse en una metáfora visible. La curva Bézier hace que el alumno
\emph{vea} la operación
$K_{IJ} \mathrel{+}= k^{(e)}_{ij}$ como una migración espacial. Las
trazas residuales acumulan un mapa que cuenta la historia del ensamblaje.

\end{moduleBox}

% ============================================================
\subsection{M6 \textbar{} Fuerzas equivalentes - Modo B (Overlay)}

\begin{moduleBox}[M6 - Fuerzas equivalentes - Modo B (Overlay)]{phaseProc}

\textbf{Concepto enseñado.} Conversión de cargas distribuidas (arista o
gravedad) en fuerzas puntuales nodales por integración con funciones
de forma.

\textbf{Patología actual.} Toplevel + dos paneles + tabla numérica que
duplica las barras. Selector de arista por combobox.

\textbf{Propuesta UX - Drag-to-load + lluvia de partículas:}
\begin{enumerate}[leftmargin=1.2em, itemsep=2pt]
  \item Sin Toplevel. Activar M6 desde el panel proc cambia el cursor
        sobre el canvas a ``modo carga''.
  \item \textbf{Drag desde fuera del elemento hacia una arista} define
        $q$ (magnitud y sentido). Al soltar:
  \item Las flechitas distribuidas se generan automáticamente sobre
        la arista; con un \emph{tween} de 1.5\,s caen como partículas
        hacia los nodos extremos, formando bolitas con tamaño
        proporcional a la fuerza integrada.
  \item Modo gravedad: switch en el panel cambia a ``lluvia uniforme''
        desde arriba, redistribuyéndose en los 9 nodos Q9 según $N_i$.
  \item \textbf{Hover sobre un nodo}: tooltip con $F_x, F_y$ exactos.
\end{enumerate}

\begin{center}
\begin{tikzpicture}[scale=0.95]
  \draw[panel] (0, 0) rectangle (15, 6.5);
  \node[text=textMuted, font=\tiny, anchor=north west]
    at (0.2, 6.4) {MeshCanvas - modo M6 activo};

  % element with edge load
  \begin{scope}[shift={(2, 1)}, scale=1.4]
    \draw[accent, fill=accent!12, line width=0.7pt]
      (0,0) rectangle (3, 2);
    \node[nodeDot] at (0,0) {}; \node[label] at (-0.2,-0.2) {1};
    \node[nodeDot] at (3,0) {}; \node[label] at (3.2,-0.2) {2};
    \node[nodeDot] at (3,2) {}; \node[label] at (3.2,2.2) {3};
    \node[nodeDot] at (0,2) {}; \node[label] at (-0.2,2.2) {4};
    \node[midDot] at (1.5,0) {};
    \node[midDot] at (3,1) {};
    \node[midDot] at (1.5,2) {};
    \node[midDot] at (0,1) {};
    \node[cnDot]  at (1.5,1) {};

    % distributed arrows on top edge
    \foreach \xa in {0.3, 0.7, 1.1, 1.5, 1.9, 2.3, 2.7} {
      \draw[selYellow, line width=0.5pt, ->]
        (\xa, 2.5) -- (\xa, 2.05);
    }
  \end{scope}

  % equivalent nodal forces (red bigger arrows)
  \begin{scope}[shift={(2, 1)}, scale=1.4]
    \draw[danger, line width=1.6pt, -{Stealth[length=4pt]}]
      (0, 2) -- (0, 3.0);
    \node[text=danger, font=\tiny, anchor=west] at (0.05, 2.65) {qL/6};

    \draw[danger, line width=2.4pt, -{Stealth[length=5pt]}]
      (1.5, 2) -- (1.5, 3.4);
    \node[text=danger, font=\tiny, anchor=west] at (1.55, 2.85) {4qL/6};

    \draw[danger, line width=1.6pt, -{Stealth[length=4pt]}]
      (3, 2) -- (3, 3.0);
    \node[text=danger, font=\tiny, anchor=west] at (3.05, 2.65) {qL/6};
  \end{scope}

  % drag indicator
  \draw[accent, dashed, line width=0.5pt, ->]
    (10, 5.5) .. controls (8, 4.5) and (6.5, 4.0) .. (5.0, 4.0);
  \node[text=accent, font=\tiny, anchor=west] at (10.1, 5.5)
    {drag desde aqui hacia la arista para aplicar q};

  % overlay legend
  \draw[panel, fill=bgDark!95, rounded corners=2pt, line width=0.4pt]
    (10.5, 1.2) rectangle (14.7, 3.0);
  \node[text=accent, font=\tiny\bfseries, anchor=north west]
    at (10.6, 2.9) {Distribucion Q9 (q cte)};
  \node[text=white, font=\tiny\ttfamily, anchor=north west]
    at (10.6, 2.55) {N1: F = qL/6};
  \node[text=white, font=\tiny\ttfamily, anchor=north west]
    at (10.6, 2.30) {N5: F = 4qL/6 (medio)};
  \node[text=white, font=\tiny\ttfamily, anchor=north west]
    at (10.6, 2.05) {N2: F = qL/6};
  \node[text=textMuted, font=\tiny, anchor=north west]
    at (10.6, 1.65) {Hover en nodo => Fx, Fy};

  \node[text=textMuted, font=\tiny, anchor=south] at (7.5, -0.4)
    {Fig. 10 - M6: drag-to-load sobre arista; particulas caen a los nodos};
\end{tikzpicture}
\end{center}

\textbf{Justificación pedagógica.} El concepto de \emph{fuerzas
equivalentes} suele explicarse como integral abstracta. Aquí el alumno
\emph{ve} la integral: las flechitas distribuidas se acumulan en cada
nodo según su peso $N_i$. Para Q9 con $q$ constante el ratio
$1\!:\!4\!:\!1$ aparece visualmente: la flecha del nodo medio es 4 veces
más grande que la de los extremos. Eso \emph{es} el resultado clásico
$L/6, 4L/6, L/6$.

\end{moduleBox}

% ============================================================
\subsection{M7 \textbar{} Discontinuidad $\sigma$ - Modo A (Toplevel 3D)}

\begin{moduleBox}[M7 - Discontinuidad sigma - Modo A (Toplevel)]{phasePost}

\textbf{Concepto enseñado.} El FEM garantiza continuidad $C^0$ de
desplazamientos pero NO de tensiones, por eso aparecen ``saltos'' en
nodos compartidos. El smoothing nodal (promediado) lo corrige a costa
de información: los saltos crudos eran un \emph{indicador cualitativo
de error}.

\textbf{Patología actual.} Slider de morph oculto detrás de un
``radiobutton''. Selector de componente en sidebar.

\textbf{Propuesta UX - Slider gigante como protagonista:}
\begin{enumerate}[leftmargin=1.2em, itemsep=2pt]
  \item Vista 3D ocupa toda la pantalla.
  \item \textbf{Slider gigante} en el footer (ancho completo) etiquetado
        ``Crudo $\leftarrow$$\rightarrow$ Promediado''. Drag continuo.
  \item Chips horizontales arriba para componente ($\sigma_x, \sigma_y,
        \tau_{xy}$, VM).
  \item \textbf{Detector de saltos automático}: aristas con discontinuidad
        $> 10\%$ parpadean en rojo cuando el slider está cerca de ``Crudo''
        — esos son los puntos donde el alumno debe refinar la malla.
\end{enumerate}

\begin{center}
\begin{tikzpicture}[scale=0.9]
  \draw[win] (0,0) rectangle (16, 7);

  % chips top
  \foreach \i/\name/\sel in {0/sigma\_x/0, 1/sigma\_y/0, 2/tau\_xy/0,
                              3/VonMises/1} {
    \pgfmathsetmacro{\xx}{0.5 + 1.6*\i}
    \ifthenelse{\equal{\sel}{1}}{
      \draw[btnPhase, fill=phasePost!40] (\xx, 6.3) rectangle ({\xx + 1.4}, 6.7);
    }{
      \draw[panel] (\xx, 6.3) rectangle ({\xx + 1.4}, 6.7);
    }
    \node[text=white, font=\small] at ({\xx + 0.7}, 6.5) {\name};
  }

  % 3D surface mockup
  \draw[panel] (0.3, 1.3) rectangle (15.7, 6.1);

  \begin{scope}[shift={(8, 3.7)}]
    % isometric mesh patches simulating discontinuous surface
    \foreach \i in {0,1,2} {
      \foreach \j in {0,1,2} {
        \pgfmathsetmacro{\h}{0.3 + 0.5*\i + 0.2*\j + 0.4*((\i==1) ? 1 : 0)}
        \pgfmathsetmacro{\xx}{-3 + 1.8*\i + 0.6*\j}
        \pgfmathsetmacro{\yy}{-1.5 + 1.0*\j - 0.3*\i}
        \fill[red!50!orange, opacity={0.5+0.1*\h}]
          (\xx, \yy) -- ({\xx + 1.6}, {\yy - 0.3}) --
          ({\xx + 1.9}, {\yy + \h}) -- ({\xx + 0.3}, {\yy + \h + 0.3}) -- cycle;
        \draw[white, line width=0.4pt]
          (\xx, \yy) -- ({\xx + 1.6}, {\yy - 0.3}) --
          ({\xx + 1.9}, {\yy + \h}) -- ({\xx + 0.3}, {\yy + \h + 0.3}) -- cycle;
      }
    }
    % blinking discontinuity edges
    \draw[danger, line width=2pt, dashed]
      (-1.2, -0.5) -- (-0.9, 0.6);
    \draw[danger, line width=2pt, dashed]
      (1.5, -0.8) -- (1.8, 0.3);
    \node[text=danger, font=\tiny\bfseries] at (-0.5, 1.5)
      {salto > 10\%};
  \end{scope}

  % slider gigante en footer
  \draw[panel] (0.3, 0.3) rectangle (15.7, 1.1);
  \node[text=accent, font=\small\bfseries, anchor=west]
    at (0.4, 0.92) {Crudo};
  \node[text=phasePost, font=\small\bfseries, anchor=east]
    at (15.6, 0.92) {Promediado};

  \draw[softGray, line width=4pt] (1.5, 0.6) -- (14.5, 0.6);
  \draw[accent, line width=4pt] (1.5, 0.6) -- (5.0, 0.6);
  \fill[selYellow, draw=white, line width=0.5pt]
    (5.0, 0.6) circle (0.18);
  \node[text=white, font=\tiny] at (5.0, 0.30) {t = 0.27};

  \node[text=textMuted, font=\tiny, anchor=south] at (8, -0.4)
    {Fig. 11 - M7: slider gigante; aristas rojas pulsantes marcan saltos};
\end{tikzpicture}
\end{center}

\textbf{Justificación pedagógica.} Mostrar el morph como slider continuo
hace que el alumno juegue con $t \in [0, 1]$ y vea las superficies
deformarse en tiempo real. Las aristas rojas pulsantes son el aporte
nuevo: convierten una visualización pasiva en un detector de errores que
sugiere acciones (``refiná aquí'').

\end{moduleBox}

% ============================================================
\subsection{M8 \textbar{} Direcciones principales + Mohr - Modo B}

\begin{moduleBox}[M8 - Direcciones principales sigma1/sigma2 - Modo B (Overlay)]{phasePost}

\textbf{Concepto enseñado.} Esfuerzos principales $\sigma_1, \sigma_2$
como autovalores del tensor, sus direcciones (ortogonales) y su
representación geométrica en el círculo de Mohr.

\textbf{Patología actual.} Toplevel separado del canvas. La tarjeta
``Estado en el nodo'' duplica los valores que ya están en el círculo.

\textbf{Propuesta UX - Vinculación bidireccional canvas $\leftrightarrow$ Mohr:}
\begin{enumerate}[leftmargin=1.2em, itemsep=2pt]
  \item Sin Toplevel. Las cruces principales se dibujan \emph{sobre el
        MeshCanvas real} en cada nodo.
  \item Overlay flotante con el círculo de Mohr del nodo seleccionado
        (arrastrable, cerrable).
  \item \textbf{Doble vínculo bidireccional}: el alumno arrastra un punto
        sobre el círculo (cualquier ángulo $2\theta$), y en la malla todas
        las cruces giran consistentemente para mostrar el estado en ese
        plano rotado.
\end{enumerate}

\begin{center}
\begin{tikzpicture}[scale=0.85]
  \draw[panel] (0, 0) rectangle (16, 7);
  \node[text=textMuted, font=\tiny, anchor=north west]
    at (0.2, 6.85) {MeshCanvas con cruces principales sigma1 (rojo) / sigma2 (azul)};

  % Mesh with crosses at nodes
  \begin{scope}[shift={(0.7, 0.7)}, scale=0.85]
    \foreach \x in {0,1,2,3,4} {
      \foreach \y in {0,1,2,3} {
        \draw[softGray!60, line width=0.3pt]
          (\x,\y) rectangle ({\x+1},{\y+1});
      }
    }
    \foreach \xx/\yy/\ang/\mA/\mB in {
      0/0/0/0.5/0.4,    1/0/-15/0.7/0.5,  2/0/-25/0.9/0.6, 3/0/-30/1.0/0.7, 4/0/-30/0.95/0.7,
      0/1/-10/0.6/0.5,  1/1/-20/0.85/0.5, 2/1/-30/1.1/0.7, 3/1/-35/1.2/0.8, 4/1/-35/1.0/0.7,
      0/2/-15/0.65/0.5, 1/2/-25/0.9/0.6,  2/2/-35/1.15/0.8,3/2/-40/1.25/0.85,4/2/-35/1.05/0.75,
      0/3/-20/0.7/0.5,  1/3/-30/0.95/0.65,2/3/-40/1.2/0.85,3/3/-45/1.3/0.95, 4/3/-40/1.1/0.8} {
      \draw[danger, line width=1pt]
        ({\xx - \mA*0.3*cos(\ang)}, {\yy - \mA*0.3*sin(\ang)}) --
        ({\xx + \mA*0.3*cos(\ang)}, {\yy + \mA*0.3*sin(\ang)});
      \draw[accent, line width=1pt]
        ({\xx - \mB*0.3*sin(\ang)}, {\yy + \mB*0.3*cos(\ang)}) --
        ({\xx + \mB*0.3*sin(\ang)}, {\yy - \mB*0.3*cos(\ang)});
      \node[circle, fill=textMuted, draw=white,
            minimum size=2.4pt, inner sep=0pt] at (\xx, \yy) {};
    }
    % selected node
    \node[selRing] at (3, 3) {};
    \node[text=selYellow, font=\tiny\bfseries] at (3.4, 3.3) {N18};
  \end{scope}

  % Mohr circle overlay top right (draggable)
  \draw[panel, fill=bgDark!95, rounded corners=2pt, line width=0.7pt,
        drop shadow={shadow xshift=1pt, shadow yshift=-1pt}]
    (9, 1.0) rectangle (15.5, 6.5);
  \node[text=accent, font=\small\bfseries, anchor=north west]
    at (9.15, 6.4) {Circulo de Mohr - N18};
  \node[text=textMuted, font=\tiny, anchor=north east]
    at (15.4, 6.4) {(arrastrable)};

  \begin{scope}[shift={(12.25, 3.7)}, scale=1.0]
    \draw[softGray, line width=0.3pt] (-2.5, 0) -- (2.5, 0);
    \draw[softGray, line width=0.3pt] (0, -1.8) -- (0, 1.8);
    \node[text=textMuted, font=\tiny, anchor=west] at (2.4, 0.15) {sigma};
    \node[text=textMuted, font=\tiny, anchor=south] at (0.15, 1.7) {tau};

    \draw[accent, line width=1pt] (0.5, 0) circle (1.5);

    % key points
    \node[circle, fill=warn, draw=white, minimum size=4pt, inner sep=0pt]
      at (1.2, 0.9) {};
    \node[text=warn, font=\tiny, anchor=west] at (1.3, 1.1) {(sx, txy)};

    \node[circle, fill=warn, draw=white, minimum size=4pt, inner sep=0pt]
      at (-0.2, -0.9) {};

    \node[circle, fill=danger, draw=white, minimum size=4pt, inner sep=0pt]
      at (2.0, 0) {};
    \node[text=danger, font=\tiny, anchor=south] at (2.0, 0.1) {sigma1};

    \node[circle, fill=accent, draw=white, minimum size=4pt, inner sep=0pt]
      at (-1.0, 0) {};
    \node[text=accent, font=\tiny, anchor=south] at (-1.0, 0.1) {sigma2};

    % draggable angle pointer
    \node[circle, fill=selYellow, draw=white, line width=0.6pt,
          minimum size=8pt, inner sep=0pt] at (0.5, 1.5) {};
    \draw[selYellow, dashed, line width=0.4pt]
      (0.5, 0) -- (0.5, 1.5);
    \node[text=selYellow, font=\tiny, anchor=west] at (0.65, 1.5)
      {2 theta};
  \end{scope}

  \node[text=textMuted, font=\tiny, anchor=south] at (8, -0.4)
    {Fig. 12 - M8: arrastrar punto en Mohr rota TODAS las cruces de la malla a la vez};
\end{tikzpicture}
\end{center}

\textbf{Justificación pedagógica.} El círculo de Mohr es la
\emph{geometría} del tensor. Convertirlo en un control activo (arrastrable)
y verlo coordinar con todas las cruces de la malla muestra que un solo
ángulo $\theta$ define el estado completo en cada punto. Pedagogía pura.

\end{moduleBox}

% ============================================================
\subsection{M9 \textbar{} Q4 vs Q9 sandbox - Modo A (Toplevel)}

\begin{moduleBox}[M9 - Q4 vs Q9 - Modo A (Toplevel)]{phasePost}

\textbf{Concepto enseñado.} Convergencia: $h$-refinement vs $p$-refinement.
Q9 converge más rápido que Q4 con menos GDLs en problemas con flexión.

\textbf{Patología actual.} El layout actual es razonable. Falta el plot
de convergencia que cierra el círculo pedagógico.

\textbf{Propuesta UX - Plot de convergencia + slider continuo:}
\begin{enumerate}[leftmargin=1.2em, itemsep=2pt]
  \item Mantener los dos paneles lado a lado.
  \item Reemplazar radiobuttons $N \in \{0, 1, 2\}$ por un \textbf{slider
        continuo} con transición animada entre niveles.
  \item Agregar un \textbf{tercer panel debajo}: gráfica log-log de
        convergencia (eje Y: $|u|_{max}$ relativo, eje X: número de GDLs).
        Para cada $N$ se acumulan dos puntos (Q4, Q9) sobre el mismo eje.
        El alumno ve la pendiente: $O(h^2)$ para Q4, $O(h^4)$ para Q9.
  \item Botón ``Tomar como referencia'': fija un resultado y los siguientes
        se miden contra él como error \%.
\end{enumerate}

\begin{center}
\begin{tikzpicture}[scale=0.9]
  \draw[win] (0,0) rectangle (16, 7);

  % Q4 panel
  \draw[panel] (0.3, 3.5) rectangle (7.8, 6.7);
  \node[text=accent, font=\small\bfseries, anchor=north west]
    at (0.4, 6.6) {Q4 - 24 elem - 60 GDL};
  \begin{scope}[shift={(4.0, 5.0)}, scale=0.5]
    \foreach \x in {0,1,2,3,4,5} {
      \foreach \y in {0,1,2,3} {
        \pgfmathtruncatemacro{\valint}{30 + 60*((\x>=2 && \x<=4 && \y==1)? 1 : 0.5)}
        \fill[red!\valint!yellow]
          (\x, \y) rectangle ({\x+1}, {\y+1});
        \draw[white, line width=0.2pt]
          (\x, \y) rectangle ({\x+1}, {\y+1});
      }
    }
  \end{scope}

  % Q9 panel
  \draw[panel] (8.2, 3.5) rectangle (15.7, 6.7);
  \node[text=phasePost, font=\small\bfseries, anchor=north west]
    at (8.3, 6.6) {Q9 - 24 elem - 200 GDL};
  \begin{scope}[shift={(11.95, 5.0)}, scale=0.5]
    \foreach \x in {0,1,2,3,4,5} {
      \foreach \y in {0,1,2,3} {
        \pgfmathtruncatemacro{\valint}{40 + 50*((\x>=2 && \x<=4 && \y==1)? 1 : 0.5)}
        \fill[red!\valint!yellow]
          (\x, \y) rectangle ({\x+1}, {\y+1});
        \draw[white, line width=0.2pt]
          (\x, \y) rectangle ({\x+1}, {\y+1});
      }
    }
  \end{scope}

  % slider
  \draw[panel] (0.3, 2.6) rectangle (15.7, 3.3);
  \node[text=accent, font=\tiny\bfseries, anchor=west]
    at (0.4, 2.95) {Refinamiento N};
  \draw[softGray, line width=4pt] (2.5, 2.95) -- (15.0, 2.95);
  \draw[accent, line width=4pt] (2.5, 2.95) -- (6.5, 2.95);
  \fill[selYellow, draw=white, line width=0.5pt]
    (6.5, 2.95) circle (0.16);
  \foreach \x/\lab in {2.5/N=0, 6.5/N=1, 10.5/N=2, 14.5/N=3} {
    \node[text=textMuted, font=\tiny] at (\x, 2.7) {\lab};
  }

  % convergence plot
  \draw[panel] (0.3, 0.3) rectangle (15.7, 2.4);
  \node[text=accent, font=\tiny\bfseries, anchor=north west]
    at (0.4, 2.3) {Convergencia log-log: |u|\_max vs GDLs};

  \begin{scope}[shift={(0.5, 0.5)}]
    \draw[softGray, line width=0.4pt] (0, 0) -- (14.5, 0);
    \draw[softGray, line width=0.4pt] (0, 0) -- (0, 1.5);
    \node[text=textMuted, font=\tiny] at (-0.05, 1.5) {error};
    \node[text=textMuted, font=\tiny] at (14.5, -0.15) {GDLs (log)};

    % Q4 curve - shallow
    \draw[accent, line width=1.2pt]
      (1, 1.2) -- (4, 0.85) -- (7, 0.55) -- (10, 0.35);
    \foreach \pt in {(1,1.2),(4,0.85),(7,0.55),(10,0.35)} {
      \node[circle, fill=accent, draw=white, minimum size=4pt, inner sep=0pt] at \pt {};
    }
    \node[text=accent, font=\tiny, anchor=west] at (10.2, 0.35) {Q4 (h\^{}2)};

    % Q9 curve - steep
    \draw[phasePost, line width=1.2pt]
      (1.5, 1.0) -- (4.5, 0.5) -- (8, 0.18) -- (11, 0.08);
    \foreach \pt in {(1.5,1.0),(4.5,0.5),(8,0.18),(11,0.08)} {
      \node[circle, fill=phasePost, draw=white, minimum size=4pt, inner sep=0pt] at \pt {};
    }
    \node[text=phasePost, font=\tiny, anchor=west] at (11.2, 0.08) {Q9 (h\^{}4)};
  \end{scope}

  \node[text=textMuted, font=\tiny, anchor=south] at (8, -0.4)
    {Fig. 13 - M9: paneles Q4/Q9 + slider continuo + plot de convergencia que cierra el ciclo};
\end{tikzpicture}
\end{center}

\textbf{Justificación pedagógica.} La pendiente del log-log \emph{es} el
orden de convergencia. Sin ese plot el alumno cree que ``Q9 es mejor'' por
intuición; con el plot lo \emph{demuestra} él mismo. Es el módulo donde
EduFEM se vuelve un laboratorio de teoría de aproximación.

\end{moduleBox}

\newpage

% ============================================================
\section{Librerías alternativas a Matplotlib}

Matplotlib es el caballo de batalla actual y debe seguir siéndolo en los
casos donde brilla: plots estáticos, LaTeX vía mathtext, gráficos
publicables. Pero hay tres áreas donde matplotlib limita la propuesta:

\begin{enumerate}[leftmargin=1.2em, itemsep=2pt]
  \item Drag de nodos a 60\,fps (M0, M6, M8).
  \item 3D rotable suave con shaders (M7, M9).
  \item Tooltips ricos sobre celdas / aristas (M5).
\end{enumerate}

\subsection{Tabla comparativa}

\begin{center}
\renewcommand{\arraystretch}{1.3}
\begin{tabular}{@{}p{3.4cm} p{4.5cm} p{3.5cm} p{3.5cm}@{}}
\toprule
\textbf{Necesidad} & \textbf{Librería} & \textbf{Pros} & \textbf{Contras} \\
\midrule
Drag fluido nodos &
\texttt{pyqtgraph} &
GPU-accel, 60\,fps, sliders sin lag, API simple &
~10\,MB (PyQt6), no integra nativo con \texttt{ttkbootstrap} \\
\midrule
3D shading + isolíneas &
\texttt{PyVista} (sobre VTK) &
Calidad ParaView, smooth shading, isovaluas 3D &
~80\,MB (VTK pesado en instalador) \\
\midrule
3D ligero rápido &
\texttt{vispy} &
OpenGL nativo, \textasciitilde 15\,MB &
API más cruda, menos ejemplos \\
\midrule
Animaciones explicativas &
\texttt{manim Community} &
Pre-render WebP, infra ya existe (\texttt{WebpPlayer}). 0\,MB runtime &
Solo offline, no interactivo \\
\midrule
Plots interactivos hover &
\texttt{plotly} + \texttt{pywebview} &
Tooltips por celda, zoom infinito &
~30\,MB, abre browser embebido \\
\midrule
Imgui-style controles &
\texttt{dearpygui} &
Look ``CAD profesional'' &
Refactor enorme, lock-in \\
\bottomrule
\end{tabular}
\end{center}

\subsection{Recomendación pragmática}

\begin{ideaBox}[Estrategia híbrida]
Mantener \textbf{matplotlib} como motor por defecto. Añadir
\textbf{pyqtgraph} como dependencia \emph{opcional} (patrón
\texttt{\_require\_pyqtgraph()} igual a \texttt{ezdxf}). Usarlo
\emph{solo} donde el drag o el slider en tiempo real lo justifica:
M0 (drag de nodo distorsionando), M6 (drag-to-load), M8 (arrastre del
punto en Mohr). Mantener matplotlib para todo lo demás. PyVista solo
si M7 escala mal con matplotlib 3D.

Para el demo-mode al abrir cada módulo: \textbf{pre-renderizar WebPs con
manim} -- encaja con la infra existente del pipeline
\texttt{tools/render\_q4q9/} y NO suma peso al runtime.
\end{ideaBox}

% ============================================================
\section{Roadmap incremental}

Migrar los nueve módulos a la vez sería temerario. Se propone un
roadmap por riesgo creciente:

\begin{center}
\begin{tikzpicture}[scale=1]
  \draw[accent, line width=3pt, ->] (0, 0) -- (15, 0);

  \foreach \x/\name/\risk/\col in {1.0/Fase 1/bajo/good,
                                    5.0/Fase 2/medio/warn,
                                    9.0/Fase 3/alto/danger,
                                    13.0/Fase 4/alto+/danger} {
    \fill[\col] (\x, 0) circle (0.2);
    \node[text=\col, font=\small\bfseries, anchor=south] at (\x, 0.3) {\name};
    \node[text=textMuted, font=\tiny] at (\x, -0.4) {riesgo \risk};
  }
\end{tikzpicture}
\end{center}

\subsection*{Fase 1 - Riesgo bajo: M7 + M8}

Postproceso, layouts simples de migrar al modo overlay. Validan el
patrón de operación desde el panel lateral sin tocar el solver ni la
arquitectura del base.

\subsection*{Fase 2 - Riesgo medio: M0 + M3}

M0 con \texttt{pyqtgraph} prueba el drag fluido. M3 prueba el modo
embed en panel lateral. Ambos son módulos pequeños donde el costo de
equivocarse es bajo.

\subsection*{Fase 3 - Riesgo alto: M1, M2, M5, M6, M9}

Una vez validados los modos en módulos chicos, se migran los más
complejos. M5 con la animación de vuelo Bézier es el más vistoso y el
que más impacto pedagógico tiene.

\subsection*{Fase 4 - Último: M4}

El integrando simbólico Q9 ya tiene problemas de renderizado LaTeX.
Migrar a la cinta transportadora requiere repensar la animación entera.

% ============================================================
\section*{Cierre}
\addcontentsline{toc}{section}{Cierre}

EduFEM tiene un núcleo FEM sólido y un ecosistema gráfico maduro. La
oportunidad pedagógica está en romper la uniformidad del layout actual
y en aprovechar el \texttt{MeshCanvas} compartido como protagonista —
los módulos no deben ser \emph{ventanas alternativas} a la GUI principal,
sino \emph{capas educativas} que se posan sobre ella. Cuando un módulo
se vuelve \emph{parte del flujo natural}, deja de ser un add-on y se
vuelve la GUI misma.

\vspace{1cm}
\begin{center}
\textit{``El mejor diseño es el que se siente inevitable''}
\end{center}

\end{document}
